% Revize 2026-10-28: obsah, matematika a návaznost; podrobnosti a časy
% viz planovani/REVIZE_PRIPRAV_1_ROCNIK.md.
% Aktuální porovnání se skutečným obsahem místních PDF Realisticky.cz:
% externi_zdroje/realisticky/01_Zakladni_poznatky/08_Mnohocleny/010805_Deleni_mnohoclenu/010805_Lekce.pdf
% Revize 2026-10-27: cílené doplnění; nové úlohy a řešení jsou vlastní tvorba.
% Metodický vzor: 4_rocnik/pripravy/Uvod_do_kombinatoriky.pdf (samostatný dokument).
% Hodina 47, 1. ročník SPŠ. Rozsah: Rozpis_vyuky_MAT_1_rocnik.tex, řádek 47.

% Plán: 1_rocnik/tematicke_plany/SPS.pdf; společná projektová šablona.

% Učebnice: ucebnice/01_Zakladni_poznatky/Algebraicke_vyrazy/Algebraické výrazy.tex

% Původ úloh: P1: přesný výraz z kostry učitele, vlastní řešení. Ostatní vlastní.

% Metodický zdroj: externi_zdroje/realisticky/01_Zakladni_poznatky/08_Mnohocleny/010805_Deleni_mnohoclenu/010805_Lekce.pdf

% URL: http://www.realisticky.cz/ucebnice/01%20Matematika%20S%C5%A0/01%20Z%C3%A1kladn%C3%AD%20poznatky/08%20Mnoho%C4%8Dleny/05%20D%C4%9Blen%C3%AD%20mnoho%C4%8Dlen%C5%AF.pdf

\documentclass[11pt,a4paper]{article}

% Spolecny styl pro ucitelske pripravy MAT (Work/MAT)
% Verze 1.1, 2026-10-23
\NeedsTeXFormat{LaTeX2e}
\ProvidesPackage{priprava}[2026/10/23 v1.1 Sablona ucitelskych priprav MAT]

% Zakladni typografie a cestina
\RequirePackage[T1]{fontenc}
\RequirePackage[utf8]{inputenc}
\RequirePackage[czech]{babel}
\RequirePackage{lmodern}
\RequirePackage{microtype}
\RequirePackage{amsmath,amssymb}
\RequirePackage{xcolor}
\RequirePackage{enumitem}
\RequirePackage[most]{tcolorbox}
\RequirePackage{tabularx,booktabs,array}
\RequirePackage{fancyhdr}
\RequirePackage{needspace}
\RequirePackage{geometry}
\RequirePackage[hidelinks]{hyperref}

\geometry{a4paper,top=17mm,bottom=18mm,left=23mm,right=23mm}
\definecolor{notegray}{gray}{0.39}
\setlength{\parindent}{0pt}
\setlength{\parskip}{2.5pt plus 1pt minus 0.5pt}
\setlist[itemize]{topsep=2pt,itemsep=1pt,parsep=0pt,leftmargin=1.45em}
\setlist[enumerate]{topsep=2pt,itemsep=2pt,parsep=0pt,leftmargin=1.65em}
\setcounter{secnumdepth}{0}
\renewcommand{\arraystretch}{1.12}

% Zahlavi: vlevo rocnik | globalni tema, vpravo cislo: nazev hodiny
% Zapati zustava zcela prazdne.
\pagestyle{fancy}
\fancyhf{}
\renewcommand{\headrulewidth}{0.35pt}
\renewcommand{\footrulewidth}{0pt}
\setlength{\headheight}{14pt}
\newcommand{\nastavPripravu}[4]{%
  \fancyhead[L]{\small #1\enspace\textbar\enspace #2}%
  \fancyhead[R]{\small #3: #4}%
}

% Minimalisticke cernobile boxy. Kazda definice zvlast.
\tcbset{commonlesson/.style={enhanced,breakable=false,colback=white,colframe=black,
  colbacktitle=white,coltitle=black,boxrule=0.45pt,
  arc=0.8mm,outer arc=0.8mm,left=3mm,right=3mm,top=1.6mm,bottom=1.6mm,
  toptitle=1mm,bottomtitle=0.7mm,fonttitle=\bfseries,
  attach boxed title to top left={xshift=2.5mm,yshift=-1.15mm},
  boxed title style={colback=white,colframe=white,boxrule=0pt,left=1mm,right=1mm}}}
\newtcolorbox{zadaniBox}[1][]{commonlesson,title={Zadání},#1}
\newtcolorbox{prikladBox}[1][]{commonlesson,title={Příklad},#1}
\newtcolorbox{definiceBox}[1][]{commonlesson,title={Definice},#1}

% Spolecny box pro zadani a jeho reseni, bez mezery mezi castmi.
% Pouziti:
%   \begin{prikladReseniBox}
%     Zadani ulohy.
%     \tcblower
%     \begin{align*} ... \end{align*}
%   \end{prikladReseniBox}
% Spodni cast ma jemny sedy podklad a delici linku; lze v ni sazet
% vice radku matematiky. Obe casti maji jen jeden spolecny ramecek.
\tcbset{solutionlesson/.style={
  commonlesson,
  enhanced jigsaw,
  breakable,
  colbacklower=black!3,
  segmentation style={solid,draw=black!25,line width=0.35pt},
  fontlower=\small\color{notegray},
  before lower={\noindent\textbf{Řešení:}\enspace},
  leftlower=3mm,rightlower=3mm
}}
\newtcolorbox{prikladReseniBox}[1][]{solutionlesson,title={Příklad},#1}
\newtcolorbox{zadaniReseniBox}[1][]{solutionlesson,title={Zadání},#1}

% Pomucky pro autora pripravy: sedive poznamky; \reseni pro starsi pripravy
\newcommand{\poznamka}[1]{\par\begingroup\small\color{notegray}\itshape
  \textnormal{\textbf{Poznámka pro učitele.}}\enspace #1\par\endgroup}
\newcommand{\reseni}[1]{\par\begingroup\small\color{notegray}
  \textbf{Očekávané řešení:}\enspace #1\par\endgroup}
\newcommand{\blok}[1]{\needspace{5\baselineskip}\subsection*{#1}}
\newcommand{\vyrok}[1]{\uv{#1}}


\nastavPripravu{1. ročník}{Algebraické výrazy}{47}{Dělení mnohočlenů se zbytkem}

\begin{document}

\blok{Výklad a souvislosti}

Při dělení nemusí mezivýsledek nakonec zmizet. Jestliže má už nižší stupeň než dělitel, další mnohočlenný člen podílu nedoplníme. Zbytek zapíšeme tak, aby šlo výsledek ověřit násobením a přičtením.

\begin{definiceBox}[title={Definice},breakable=false]
\textbf{Dělení se zbytkem} zapisujeme $A=BQ+R$, kde $B$ není nulový mnohočlen a zbytek $R$ je nulový, nebo má stupeň menší než $B$. Mnohočlen $Q$ je podíl, $R$ zbytek.
\end{definiceBox}


Z identity plyne $A(x)/B(x)=Q(x)+R(x)/B(x)$ pro $B(x)\ne0$. Při lineárním děliteli je nenulový zbytek konstantní. Samotné $Q$ není celou hodnotou původního zlomku, pokud $R\ne0$.

\blok{Řešené příklady}

\begin{prikladReseniBox}[title={Příklad 1},breakable=false]
Vyděl $(p^3+2p^2-2p-1):(p+1)$.
\tcblower
Písemné schéma:
\[
\begin{array}{r r r r l}
 p^3&+2p^2&-2p&-1&:\,(p+1)=p^2+p-3\quad\text{(zbytek }2\text{)}\\
 -\bigl(p^3&+p^2&&\bigr)&\text{odečítáme }p^2(p+1)\\ \hline
 &p^2&-2p&-1&\\
 &-\bigl(p^2&+p&\bigr)&\text{odečítáme }p(p+1)\\ \hline
 &&-3p&-1&\\
 &&-\bigl(-3p&-3\bigr)&\text{odečítáme }(-3)(p+1)\\ \hline
 &&&2&
\end{array}
\]
Nejvyšší členy dělíme: $p^3:p=p^2$, $p^2:p=p$, $-3p:p=-3$.
Po každém kroku odečítáme celý součin; například $-1-(-3)=2$.
Zbytek $2$ má stupeň $0<1$, proto končíme. Stejné kroky jako rovnice:
\begin{align*}
(p^3+2p^2-2p-1)-p^2(p+1)&=p^2-2p-1,\\
(p^2-2p-1)-p(p+1)&=-3p-1,\\
(-3p-1)-(-3)(p+1)&=2.
\end{align*}
Tedy $Q=p^2+p-3$, $R=2$ a původní podíl je $p^2+p-3+\frac2{p+1}$ při $p\ne-1$. Kontrola: $(p+1)(p^2+p-3)+2=p^3+2p^2-2p-1$.
\end{prikladReseniBox}


\begin{prikladReseniBox}[title={Příklad 2},breakable=false]
Vyděl $(x^3+2x+1):(x^2+1)$.
\tcblower
První člen podílu je $x$. Po odečtení $x(x^2+1)$ zbude $x+1$, jehož stupeň je $1<2$. Proto $Q=x$, $R=x+1$. Jmenovatel $x^2+1$ je pro všechna reálná $x$ kladný.
\end{prikladReseniBox}


\blok{Samostatné procvičení}

\begin{zadaniReseniBox}[title={Procvičení},breakable=false]
\begin{enumerate}[label=\alph*)]
\item Vyděl $(x^2+1):(x-1)$.
\item Vyděl $(2x^2+3x+4):(x+1)$.
\item Vyděl $(x^3-1):(x^2+1)$.
\end{enumerate}
\tcblower
\begin{enumerate}[label=\alph*)]
\item $Q=x+1$, $R=2$, protože $(x-1)(x+1)+2=x^2+1$. Pro podíl hodnot je $x\ne1$.
\item $Q=2x+1$, $R=3$; $(x+1)(2x+1)+3=2x^2+3x+4$. Pro zlomek $x\ne-1$.
\item $Q=x$, $R=-x-1$. Kontrola: $x(x^2+1)-x-1=x^3-1$.
\end{enumerate}
\end{zadaniReseniBox}


\poznamka{Zbytek nemusí být číslo, jak ukazuje kvadratický dělitel. Vždy ověř podmínku stupně i identitu $A=BQ+R$. U P1 lze navíc dosadit $p=-1$ do této identity: dělenec má hodnotu $2$, stejně jako zbytek; do původního zlomku však toto číslo dosadit nelze.}


\poznamka{Základ: původní P1 s podrobným schématem, P2 s kvadratickým dělitelem a procvičení a), b). Původní c) je rezerva k procvičení nečíselného zbytku. Domácí trojice střídá nenulový a nulový zbytek; vždy kontrolujeme stupeň i identitu, podmínku pro zlomek posuzujeme zvlášť.}

\blok{Domácí procvičení}

\begin{zadaniReseniBox}[title={Domácí dělení se zbytkem},breakable=false]
\begin{enumerate}[label=\alph*)]
\item Vyděl $(x^2+2x+5):(x+1)$ a ověř identitu $A=BQ+R$.
\item Vyděl $(2x^3-x+2):(x^2+1)$. Proč už nelze pokračovat?
\item Vyděl $(x^3-2x^2+1):(x-1)$; urči podmínku podílu hodnot.
\end{enumerate}
\tcblower
\begin{enumerate}[label=\alph*)]
\item $Q=x+1$, $R=4$; $(x+1)^2+4=x^2+2x+5$, $x\ne-1$.
\item $Q=2x$, $R=-3x+2$; $2x(x^2+1)-3x+2=2x^3-x+2$. Stupeň zbytku $1<2$; dělitel je na $\mathbb R$ kladný.
\item $Q=x^2-x-1$, $R=0$; $x\ne1$. Mezivýsledky jsou $-x^2+0x+1$, $-x+1$, $0$; zpětné násobení vrací dělenec.
\end{enumerate}
\end{zadaniReseniBox}

\end{document}
